\section{Background}
\paragraph{Sycophancy as an Alignment Failure} \citet{sharma2023towardsunderstanding} show  that models trained with human feedback exhibit sycophancy across free-form generation tasks and that both human raters and preference models may favor responses that match user views over factual ones. Similar single-turn evaluations such as BullshitBench~\citep{gostev2026bullshitbench} study whether models reject nonsensical or false-premise prompts observing if the models over-defer to the user's framing rather than challenging it. While this line of work establishes sycophancy as a real alignment failure, they do not explain when a model is most vulnerable to abandoning a correct answer during an interaction.

\paragraph{Multi-Turn Sycophancy and Pressure-Based Evaluation}
Sycophancy goes beyond single-turn factuality failure. TRUTH DECAY~\citep{liu2025truthdecayquantifyingmultiturn} evaluates sycophancy in extended dialogues where models face iterative user feedback, challenges, and persuasion. SYCON Bench~\citep{hong2025syconbench} further probes multi-turn sycophancy using flip-based metrics such as Turn of Flip and Number of Flips, measuring how quickly and how often models conform to sustained user pressure. SWAY~\citep{bhalla2026sway} provides a counterfactual linguistic measure of sycophancy by comparing model agreement under positive versus negative framing, isolating pressure-induced agreement. These works establish that models can drift or flip under pressure. Our work investigates whether such flips are predictable from the model's uncertainty on a specific question and reasoning trajectory across escalating pressure from the user.

\paragraph{Belief Updating Under Persuasion and Feedback} Prior work \citet{xu2023earthisflat} show that models' correct factual beliefs can be manipulated through multi-turn persuasive misinformation. In contrast, \citet{jiang2025feedbackfriction} study that models often fail to incorporate high-quality corrective feedback when the model is confident, even when the feedback is backed by ground-truth answers. Together, these results suggest that LLM belief updating is miscalibrated in both directions and that models may resist valid correction while remaining vulnerable to invalid persuasion. These findings motivate our focus on how pressure and uncertainty shape whether models revise their answers.

\paragraph{Uncertainty Calibration in Language Models}  A separate line of work studies whether language models can express calibrated uncertainty. \citet{tian2023justask} show that verbalized confidence can be better calibrated than conditional probabilities for RLHF-tuned models. ConfTuner~\citep{li2025conftuner} improves verbalized confidence through a tokenized Brier-score objective, while Uncertainty Distillation~\citep{hager2025uncertaintydistillationteachinglanguage} trains models to express semantic confidence derived from repeated sampling and semantic normalization. These works motivate our use of uncertainty as a black-box behavioral signal but they primarily treat uncertainty as a calibration target to make confidence match correctness. Our work investigates uncertainty as a signal of vulnerability to social pressure.

\paragraph{Reasoning-Trace and Process-Level Confidence} Recent work treats uncertainty as a dynamic signal over the reasoning process. \citet{yoon2025reasoningconfidence} find that reasoning models are often better calibrated than non-reasoning models and become increasingly calibrated as their chain-of-thought unfolds. However, \citet{mei2025reasoninguncertainty} show that deeper reasoning can also increase overconfidence on incorrect answers, suggesting that chain-of-thought does not automatically produce reliable uncertainty. Other work explicitly models reasoning-time confidence: Temporalizing Confidence~\citep{mao2025temporalizingconfidence} treats stepwise chain-of-thought confidence as a temporal signal, Thought Calibration~\citep{wu2025thoughtcalibration} probes partial reasoning traces for calibrated stopping, and PRM calibration work~\citep{park2025knowwhatyoudontknow} estimates prefix-level success probabilities. Complementary work trains models to emit local uncertainty
signals during reasoning~\citep{guo2026llmsshouldexpressuncertainty}. Influenced by prior work, we ask whether chain-of-thought belief trajectories reveal the path to capitulation under escalating pressure.
